\documentclass[11pt]{article}
\usepackage[margin=1in]{geometry}
\usepackage[T1]{fontenc}
\usepackage[utf8]{inputenc}
\usepackage{lmodern}
\usepackage{amsmath,amssymb,mathtools,booktabs,hyperref}

\title{Penalized Trend Estimation from First Principles: A Step-by-Step Companion}
\author{Heriberto Espino Montelongo}
\date{\today}

\newcommand{\R}{\mathbb{R}}

\begin{document}
\maketitle

\begin{abstract}
This companion document explains the mathematical and computational ideas used
in the Trend Estimation project for a reader who has not previously studied
penalized smoothing. It develops finite differences, quadratic penalties,
closed-form trend estimation, the meaning of the smoothing parameter, temporal
validation, analytic sensitivity, and numerical optimization. It also explains
why the pure penalized model and the Guerrero model with an estimated drift
must be treated separately.
\end{abstract}

\section{Glossary and notation}
Let $y=(y_1,\ldots,y_n)^\top\in\R^n$ be an observed time series and let
$t=(t_1,\ldots,t_n)^\top$ denote an unknown smooth trend.

\begin{itemize}
\item $D_d$: matrix that computes finite differences of order $d$.
\item $Q=D_d^\top D_d$: quadratic roughness-penalty matrix.
\item $\lambda\ge 0$: penalty or smoothing parameter.
\item $S_\lambda=(I+\lambda Q)^{-1}$: smoothing matrix for the pure model.
\item forecast origin: last time index whose observation is allowed to enter a fit.
\item look-ahead leakage: use of information that would not have been known at the forecast origin.
\end{itemize}

\section{Why finite differences measure curvature}
For first differences,
\[
\Delta t_i=t_{i+1}-t_i.
\]
Second differences are
\[
\Delta^2t_i=t_{i+2}-2t_{i+1}+t_i.
\]
If $\Delta^2t_i=0$ everywhere, the sequence has constant first difference and
therefore follows a discrete straight line. More generally, forcing
$\Delta^dt=0$ produces a discrete polynomial of degree at most $d-1$.

\section{The pure penalized smoother}
We balance fidelity to the observations against roughness:
\begin{equation}
J(t;\lambda)=\|y-t\|_2^2+\lambda\|D_dt\|_2^2.
\end{equation}
The first term is small when $t$ follows the data. The second term is small when
$t$ has small order-$d$ differences. Increasing $\lambda$ therefore prioritizes
smoothness more strongly.

Because
\[
\|D_dt\|_2^2=t^\top D_d^\top D_dt=t^\top Qt,
\]
we have
\[
J(t;\lambda)=(y-t)^\top(y-t)+\lambda t^\top Qt.
\]
Differentiating with respect to $t$ gives
\[
\nabla_tJ=2(t-y)+2\lambda Qt.
\]
Setting the gradient equal to zero,
\[
(I+\lambda Q)t=y.
\]
Hence
\begin{equation}
\boxed{\widehat t_\lambda=(I+\lambda Q)^{-1}y.}
\end{equation}
In code we solve the corresponding linear/spectral system; we do not explicitly
form a matrix inverse.

\section{What happens when lambda changes?}
The trend itself is a function of $\lambda$. Write
\[
S_\lambda=(I+\lambda Q)^{-1},\qquad \widehat t_\lambda=S_\lambda y.
\]
The derivative of a matrix inverse is
\[
\frac{dA^{-1}}{d\lambda}
=-A^{-1}\frac{dA}{d\lambda}A^{-1}.
\]
Since $A=I+\lambda Q$ and $A'=Q$,
\[
S_\lambda'=-S_\lambda Q S_\lambda.
\]
Therefore
\begin{equation}
\boxed{\widehat t_\lambda'=-S_\lambda Q\widehat t_\lambda.}
\end{equation}
Differentiating one more time gives
\begin{equation}
\boxed{\widehat t_\lambda''=2S_\lambda Q S_\lambda Q\widehat t_\lambda.}
\end{equation}
The implementation checks these expressions against finite-difference
approximations before they are used by Newton-type optimization.

\section{Why optimize log-lambda?}
A valid penalty satisfies $\lambda>0$. Introduce
\[
\theta=\log\lambda,\qquad \lambda=e^\theta.
\]
Then every real $\theta$ maps automatically to a positive $\lambda$. If
$g(\theta)=f(e^\theta)$,
\[
g'(\theta)=\lambda f'(\lambda),
\]
and
\[
g''(\theta)=\lambda f'(\lambda)+\lambda^2f''(\lambda).
\]
This allows Brent-type minimization or Newton iterations to work on an
unconstrained logarithmic scale while preserving positivity.

\section{Why validation must respect time}
Suppose the model is evaluated on observations $y_{T+1},\ldots,y_{T+h}$. A
valid experiment fits the model using only
\[
y_1,\ldots,y_T.
\]
The future block is then used only to compute an error. Fitting a smoother on
the entire vector first and subsequently calling the last observations
``validation'' would leak future information into the fitted trend.

Rolling-origin evaluation repeats the valid procedure at several origins:
\[
1{:}T_1\rightarrow T_1{+}1{:}T_1{+}h,
\quad
1{:}T_2\rightarrow T_2{+}1{:}T_2{+}h,
\quad\ldots
\]
This is the principal validation design for financial experiments in the
project.

\section{Why Guerrero is a different model}
The original project uses
\begin{equation}
\widehat\tau_{\lambda,d}
=(I+\lambda Q)^{-1}
\left(y+\lambda\widehat mD_d^\top\mathbf1\right).
\end{equation}
Here $\widehat m$ represents an estimated order-$d$ difference drift. The code
re-estimates this quantity from the fitted trend, so in general
\[
\widehat m=\widehat m(\lambda).
\]
Consequently, differentiating $\widehat\tau$ requires the product/chain rule
for both the inverse matrix and $\widehat m(\lambda)$. The pure-model formula
cannot simply be copied. The project therefore exposes two explicit estimators:
\texttt{PurePenalizedTrend} and \texttt{GuerreroTrend}.

\section{From smoothing to financial decisions}
The eventual research pipeline is
\[
y\longrightarrow \widehat t_\lambda
\longrightarrow \text{forecast or signal}
\longrightarrow \text{decision}.
\]
There are then several possible definitions of ``optimal'' smoothing:
\[
\lambda_{\rm GCV},\qquad
\lambda_{\rm forecast},\qquad
\lambda_{\rm decision}.
\]
A central scientific question is whether these coincide. The project does not
assume that they do.

\section{Implementation map}
The reusable mathematics belongs in \texttt{src/trend\_estimation}. Tests
verify the operators, estimators, derivatives, and chronological splits.
Experiments call the installed package. The formal paper states results
compactly; this document records the derivations and implementation logic in a
form suitable for learning and auditing.

\end{document}
